\documentclass[11pt]{paper}

\usepackage[T1]{fontenc}

\usepackage{graphicx}
\usepackage{amssymb}
\usepackage{amstext}
\usepackage{amsmath}
\usepackage{a4wide,color}
\usepackage[utf8]{inputenc}
\usepackage[frenchb]{babel}
\usepackage{xspace}
\usepackage{anysize}
\usepackage{tabularx}
\usepackage{multirow}
\usepackage{fancybox}
\usepackage{fancyhdr}
\usepackage{bbding}

\usepackage{color}
\usepackage{float}

\usepackage[toc,page]{appendix} 
\usepackage{lscape}

\usepackage{placeins}

\usepackage{listings}
\lstdefinestyle{customc}{
  belowcaptionskip=1\baselineskip,
  breaklines=true,
  frame=L,
  xleftmargin=\parindent,
  language=C,
  showstringspaces=false,
  basicstyle=\tiny\ttfamily,
  keywordstyle=\bfseries\color{green!40!black},
  commentstyle=\itshape\color{purple!40!black},
  identifierstyle=\color{blue},
  stringstyle=\color{orange},
  tabsize=2
}

\usepackage{todonotes}

\usepackage[framemethod=TikZ]{mdframed}

\lstset{% general command to set parameter(s)
	basicstyle=\footnotesize, % print whole listing small
	keywordstyle=\color{magenta}\bfseries, % underlined bold black keywords
	identifierstyle=, % nothing happens
	commentstyle=\color{black}, % white comments
	stringstyle=\ttfamily, % typewriter type for strings
	showstringspaces=false % no special string spaces
} 

\lstdefinelanguage{aald}
	{morekeywords={system, implementation},
	sensitive=false,
	morecomment=[l]{//},
	morecomment=[s]{/*}{*/},
	morestring=[b]",
}

\usepackage[colorlinks=true]{hyperref}

\usepackage[normalem]{ulem}
\usepackage{color}

\definecolor{Fond}{gray}{0.7}

\renewcommand{\floatpagefraction}{.99}
\renewcommand{\textfraction}{.01}
\newcommand{\modif}[1]{\textcolor{red}{\uline{#1}}}

\newcommand{\version}{1.0.1{ }}


\pagestyle{fancy}
\fancyhf{}
\fancyhead[RE,RO]{\thepage}
\fancyhead[LE]{}
\fancyhead[LO]{Temps réel }
\fancyfoot[LO]{Centrale Nantes -- \today}
\fancypagestyle{plain}{%
  \fancyhf{} % get rid of headers
  \renewcommand{\headrulewidth}{0pt} % and the line
}

\newenvironment{maliste}%
{ \begin{list}%
	{\ArrowBoldRightStrobe}%
	{\setlength{\labelwidth}{30pt}%
	 \setlength{\leftmargin}{35pt}%
	 \setlength{\itemsep}{\parsep}}}%
{ \end{list} }

 \setlength{\parskip}{1.2ex plus .9ex minus .4ex}

\newcommand{\ocv}{{\tt OpenCV}\xspace} 
\newcommand{\alsa}{{\tt ALSA}\xspace} 
\newcommand{\aadl}{{\tt AADL}\xspace} 
\newcommand{\xeno}{{\tt Xenomai}\xspace} 
\newcommand{\iic}{I$^2$C\xspace} 


\title{{\Large Systèmes interconnectés -- SINT }\\
\vspace{-0.4em}
{\large Sujet de TP 2 : \iic}\\
{\small Centrale Nantes}\\
\small P.-E. Hladik, \small{\texttt{pehladik@ec-nantes.fr}}\\
---\\
{\scriptsize  Version \version (\today)}\\
}


\begin{document}

\maketitle

%\begin{table}[htdp]
%\caption{Suivi des modifications}
%\begin{center}
%\begin{tabular}{|c|c|c|p{8cm}|}
%\hline
%Date & Version & Auteur & Modifications\\
%\hline
%JJ/MM/12 & 0.3.X & P.-E. Hladik&  \\
%\hline
%\end{tabular}
%\end{center}
%\label{default}
%\end{table}%

%%%%%%%%%%%%%%%%%%%%%%%%%
  \section{Objectifs pédagogiques}
%%%%%%%%%%%%%%%%%%%%%%%%%


\begin{itemize}
	\item lecture de datasheet pour configurer un composant via \iic
	\item utilisation d'une bibliothèque pour communiquer via \iic
\end{itemize}

%%%%%%%%%%%%%%%%%%%%%%%%%
  \section{Ce qu'il y a à rendre}
%%%%%%%%%%%%%%%%%%%%%%%%%

Vous déposez sur hippocampus un document pdf avec les parties pertinente de votre code commenté.

%%%%%%%%%%%%%%%%%%%%%%%%%
  \section{Ce qu'il y a à faire}
%%%%%%%%%%%%%%%%%%%%%%%%%

Ce deuxième TP s’effectue sur la carte Feather M0 Wifi. Vous pouvez utiliser le code que vous avez déjà produit dans le TP~1.

Réaliser une application qui permet d'afficher les mesures de l'accéléromètre qui est sur la carte. Vous aurez besoin du datasheet de l'accéléromètre MMA7455L qui est disponible sur le serveur pédagogique. La communication se fait avec une fréquence de 400 kHz.

Implémenter les fonctions pour :
\begin{itemize}
	\item configurer la sensibilité de l'accéléromètre (prévoir une fonction générique),
	\item passer le capteur dans le mode {\it Measurement Mode} afin de débuter les mesures,
	\item lire les valeurs sur les 3 axes,
	\item afficher vos noms sur la première ligne de l'afficheur LCD, et sur la 2e ligne l'accélération en X, suivie de l'accélération en Y et de l'accélération en Z (attention : le résultat de ces mesures est un nombre signé sur 8 bits).\\
\end{itemize}

%%%%%%%%%%%%%%%%%%%%%%%%%
  \section{Ce qu'il y a à faire (la suite)}
%%%%%%%%%%%%%%%%%%%%%%%%%

Utiliser l'analyseur logique pour observer chaque type de communication (configuration, lecture des valeurs). Réaliser des captures d'écran et mettez les dans votre rapport.

Il existe une bibliothèque Arduino pour gérer le capteur (\href{https://www.arduino.cc/reference/en/libraries/mma7455-sensor-library/}{documentation}). Chargez la bibliothèque (\href{https://downloads.arduino.cc/libraries/github.com/ricki-z/MMA7455_sensor_Library-0.0.2.zip}{https://downloads.arduino.cc/libraries/github.com/ricki-z/MMA7455\_sensor\_Library-0.0.2.zip}) et regardez le fichier d'entête proposé. Seriez-vous capable d'implémenter cette bibliothèque pour la carte Feather ? Regardez comment elle a été codé.

\end{document}